\subsection{Comparison to Manley--Bhattacharya Predictions}\label{sec:mb_comparison}
The Manley--Bhattacharya (MB) calculation combines leading-twist longitudinal amplitudes involving $\widetilde H$ and $\widetilde E$, including the charged-pion pole, with transverse amplitudes at twist~3. The transverse sector contains both the chiral-odd $H_T$ and $\bar E_T$ contributions coupled to a twist-3 pion distribution amplitude, and additional chiral-even contributions from moments of generalized transverse-momentum dependent distributions (GTMDs) coupled to a twist-2 pion distribution amplitude. In the supplied implementation, endpoint integrals are regulated with $\delta=\langle p_\perp^2\rangle/Q^2$; separate chiral-even and chiral-odd regulator parameters and a transversity $t$ slope are fitted to CLAS6 unpolarized $\pi^0$ cross sections. The displayed theory bands propagate the parameter-fit uncertainty, rather than representing the full model uncertainty.
The description above summarizes preliminary theory material supplied for the $\pi^+$ comparison; numerical neutron-channel predictions and their kinematic evaluation must be supplied and checked before any curves are shown here. The published GK framework instead describes longitudinal amplitudes with $\widetilde H$, $\widetilde E$ and the charged-pion pole, and important transverse amplitudes with chiral-odd GPDs and a twist-3 pion wave function~\cite{Goloskokov:2009ia,Goloskokov:2011rd}. The additional chiral-even GTMD-moment contributions in the MB treatment provide a distinct mechanism to examine in the transverse cross section. The comparison will retain the measured lab-axis convention, finite-bin kinematics, polarization definitions, and the distinction between parameter-fit bands and full model uncertainty. No numerical agreement or disagreement is claimed at this stage.
% Add neutron-channel predictions, figures and results after extraction.
